\documentclass[11pt,a4paper]{article}
\usepackage{amsmath,amssymb,amsthm}
\usepackage[margin=2.5cm]{geometry}
\usepackage{hyperref}

\newtheorem{definition}{Definition}[section]
\newtheorem{theorem}[definition]{Theorem}
\newtheorem{proposition}[definition]{Proposition}
\newtheorem{corollary}[definition]{Corollary}
\newtheorem{remark}[definition]{Remark}
\newtheorem{conjecture}[definition]{Conjecture}

\title{Hyperseed Formalization:\\
  God Doesn't Double-Count Evidence}
\author{ProtomegaTron (automated formalization)}
\date{2026-09-19 (deepened v2)}

\begin{document}
\maketitle

\begin{abstract}
We formalize Ben Goertzel's ``God Doesn't Double-Count Evidence'' (2022-01-15)
within the Hyperseed ontology, incorporating d-calculus extensions from
Hyperseed v2. The central claim---that rational inference must not count
the same evidence twice, and that this principle is mirrored by conservation
laws in physics---is given a fiber-geometric interpretation where evidence
independence is characterized by homotopy class membership and evidence
conservation becomes a divergence-free current condition.
\end{abstract}

\tableofcontents

\section{Evidence bundle and fiber independence}

\begin{definition}[Evidence bundle]
Let $B$ be the epistemic base space (the space of propositions or hypotheses).
An \emph{evidence bundle} is a fiber bundle $\pi : E \to B$ where each fiber
$E_b = \pi^{-1}(b)$ consists of the evidence items bearing on hypothesis $b$.
\end{definition}

\begin{definition}[Evidence independence]
Two evidence items $e_1, e_2 \in E_b$ are \emph{independent} if there exists
no derivation path connecting them within the total space $E$:
\[
  e_1 \perp e_2 \iff \nexists\, \gamma : [0,1] \to E \text{ with }
  \gamma(0) = e_1,\; \gamma(1) = e_2,\;
  \pi \circ \gamma = \text{const.}
\]
\end{definition}

\begin{proposition}[No double-counting = fiber independence --- claims 1, 3, 6, 8]
\label{prop:nodc}
Independent evidence corresponds to independent fibers:
\[
  e_1 \perp e_2 \implies \text{fuse freely}; \quad
  e_1 \not\perp e_2 \implies \text{single count (dedup required)}
\]
Double-counting is the error of treating dependent evidence items
($e_1 \not\perp e_2$) as independent, which inflates confidence
beyond what the evidence warrants.
\end{proposition}

\begin{proof}[Sketch]
If $e_1$ and $e_2$ share a derivation path $\gamma$, then they carry
overlapping information. Counting both at full weight assigns total
weight $w_1 + w_2$ when the actual information content is at most
$w_1 + w_2 - w_\gamma$ where $w_\gamma$ is the weight of the shared
path. The inflation is $w_\gamma > 0$.
\end{proof}

\section{Homotopy independence certificate}

\begin{definition}[Evidence path]
An \emph{evidence path} from source $s$ to hypothesis $b$ is a path
$\gamma : [0,1] \to E$ with $\gamma(0) \in E_s$ (source data) and
$\pi(\gamma(1)) = b$.
\end{definition}

\begin{proposition}[Homotopy criterion --- claims 4, 7, 9]
\label{prop:homotopy}
Two evidence paths $\gamma_1, \gamma_2$ that are homotopic
(connected by a continuous deformation through shared intermediate
evidence) must not be counted separately:
\[
  \gamma_1 \simeq \gamma_2 \implies \text{single count}; \quad
  \gamma_1 \not\simeq \gamma_2 \implies \text{independent contributions}
\]
The homotopy class $[\gamma]$ serves as the \emph{canonical independence
certificate}: evidence contributions are counted per homotopy class,
not per path.
\end{proposition}

\begin{remark}
This gives a principled geometric answer to ``when does evidence overlap?''
Two derivation paths overlap exactly when they are homotopic---i.e., when
one can be continuously deformed into the other through shared intermediate
lemmas. The answer is topological, not syntactic.
\end{remark}

\begin{corollary}[Evidence weight formula]
The correct total evidence weight for hypothesis $b$ from sources
$\{s_i\}$ is:
\[
  W(b) = \sum_{[\gamma] \in \pi_1(E, b)} w_{[\gamma]}
\]
where the sum runs over \emph{homotopy classes} of evidence paths,
not over individual paths. Each class contributes its weight once.
\end{corollary}

\section{Base-space conservation constraint}

\begin{proposition}[Information budget --- claims 2, 5, 10]
\label{prop:budget}
Physical laws enforce non-redundancy as a base-space constraint.
At each base point $b \in B$, the total evidence weight is bounded:
\[
  W(b) = \sum_{[\gamma]} w_{[\gamma]} \leq W_{\max}(b)
\]
where $W_{\max}(b)$ is the information capacity at base point $b$,
determined by physical constraints (channel capacity, observation
time, detector resolution, etc.).
\end{proposition}

\begin{remark}[Physics as bookkeeping]
Conservation laws (energy, momentum, charge) are reinterpreted as
the universe's mechanism for enforcing $W(b) \leq W_{\max}(b)$
within every local reference frame. ``God doesn't cook the books.''
\end{remark}

\section{d-calculus formulation (Hyperseed v2)}

\begin{definition}[Evidence connection]
An \emph{evidence connection} on the bundle $\pi : E \to B$ is a
connection $\nabla$ that specifies how evidence transforms under
parallel transport between base points (change of reference frame
or hypothesis).
\end{definition}

\begin{proposition}[Divergence-free evidence current --- claim 11]
\label{prop:divfree}
Evidence conservation is expressed as a divergence-free current
condition:
\[
  \nabla \cdot J_{\text{evidence}} = 0
\]
where $J_{\text{evidence}}$ is the evidence current (flux of
evidential weight through the base space). No evidence is
created or destroyed in transport.
\end{proposition}

\begin{proposition}[Holonomy detector --- claim 12]
\label{prop:holonomy}
Non-trivial holonomy in the evidence connection reveals
double-counting:
\[
  \operatorname{Hol}_\gamma(\nabla) \neq \mathrm{id}
  \implies \text{evidence ``leak'' along loop } \gamma
\]
Transporting evidence around a closed loop and getting a different
answer means the books are being cooked. The no-double-counting
principle demands:
\[
  \operatorname{Hol}_\gamma(\nabla) = \mathrm{id}
  \quad \forall \text{ contractible loops } \gamma
\]
within each reference frame.
\end{proposition}

\begin{proposition}[Curvature = evidence interaction --- claim 13]
\label{prop:curvature}
The curvature $F_\nabla$ of the evidence connection measures
evidence interaction:
\[
  F_\nabla = 0 \implies \text{fully independent evidence (free fusion)}
\]
\[
  F_\nabla \neq 0 \implies \text{interacting evidence (holonomy correction needed)}
\]
In the flat case, all evidence can be fused without adjustment.
In the curved case, fusion requires a holonomy correction factor
$\operatorname{Hol}_\gamma(\nabla)$ to compensate for evidence
interaction along the path.
\end{proposition}

\begin{corollary}[Corrected fusion formula]
When the evidence connection has non-zero curvature, the corrected
fusion weight for two evidence paths $\gamma_1, \gamma_2$ arriving
at the same hypothesis $b$ is:
\[
  w_{\text{fused}} = w_1 + w_2
    - \langle w_1, \operatorname{Hol}_{\gamma_1^{-1} \cdot \gamma_2}(\nabla) \cdot w_2 \rangle
\]
where the correction term accounts for the shared derivation content
detected by the holonomy along the composite path $\gamma_1^{-1} \cdot \gamma_2$.
\end{corollary}

\section{Evidence-energy duality}

\begin{conjecture}[Evidence-energy isomorphism --- claim 14]
Evidence is to logic what energy is to physics. Both are conserved
quantities whose conservation prevents systematic error:
\begin{itemize}
  \item Energy conservation prevents perpetual motion machines.
  \item Evidence conservation prevents perpetual confidence machines
    (systems that generate unbounded confidence from finite data).
\end{itemize}
The fiber bundle framework places both in the same geometric setting:
physical energy as a section of a physical bundle, logical evidence
as a section of the epistemic bundle $E \to B$, with parallel
conservation laws in each.
\end{conjecture}

\begin{remark}
This conjecture is developed into a full mathematical framework in
Goertzel's later ``Evidence Is to Logic What Energy Is to Physics''
(2026), which derives aspects of quantum-mechanical structure from
evidence conservation axioms.
\end{remark}

\end{document}
